\documentclass[11pt]{article}

\usepackage{amsmath,amsthm,amssymb}
\usepackage[margin=1in]{geometry}

\newtheorem{theorem}{Theorem}
\newtheorem{proposition}[theorem]{Proposition}
\newtheorem{corollary}[theorem]{Corollary}
\newtheorem{lemma}[theorem]{Lemma}
\theoremstyle{definition}
\newtheorem{definition}[theorem]{Definition}
\newtheorem{example}[theorem]{Example}
\newtheorem{remark}[theorem]{Remark}

\newcommand{\Hom}{\operatorname{Hom}}
\newcommand{\End}{\operatorname{End}}
\newcommand{\id}{\mathrm{id}}
\newcommand{\C}{\mathcal{C}}
\newcommand{\D}{\mathcal{D}}
\newcommand{\K}{\mathcal{K}}
\newcommand{\E}{\mathcal{E}}
\newcommand{\Set}{\mathbf{Set}}
\newcommand{\Poly}{\mathbf{Poly}}
\newcommand{\Span}{\mathbf{Span}}
\newcommand{\Cat}{\mathbf{Cat}}
\newcommand{\Sk}{\mathrm{Sk}}
\newcommand{\ob}{\operatorname{ob}}
\newcommand{\Pow}{\mathsf{P}}
\newcommand{\Dist}{\Delta}

\title{The Karamlou--Shah no-go is not an $H^2$ obstruction:\\
a level analysis, and the correct scope of the\\
cohomological closure theorem}
\author{MacBeth\\\small Kodamai}
\date{12 June 2026}

\begin{document}
\maketitle

\begin{abstract}
The companion note \cite{Gcohom} classifies the global-closure obstruction (G) to
a strict factorization $\K=\C\bowtie\D$ as a cohomology class
$[\omega]\in H^2(\Sk_\C;\D)$. A natural and tempting conjecture --- recorded as a
prove target --- asserted that the Karamlou--Shah no-go theorems (the list, stream,
and tree directed containers do \emph{not} distribute over the distribution monad
$\Dist$ or the nondeterminism monad $\Pow$) are \emph{instances} of this
obstruction: \textnormal{(G)} failures witnessed by a nontrivial $[\omega]$. We
\emph{refute} this conjecture, and we do so in a way that sharpens rather than
weakens the program. The refutation is a \emph{level theorem}. The cohomological
criterion is defined only when the prescribed right factor $\D$ is a small category
--- equivalently (Ahman--Uustalu) a directed container, equivalently a polynomial
comonad. We prove by elementary means that the finite powerset monad $\Pow$ and the
distribution monad $\Dist$ are \emph{not} polynomial functors; hence they cannot be
presented as a wide subcategory $\D$, and the crucial step ``express the monad as
$\D$'' is unsatisfiable. The Karamlou--Shah obstruction therefore lives strictly
\emph{below} the local-freeness condition \textnormal{(L)}: it is a
non-representability (failure-to-be-a-directed-container) phenomenon in the wrong
bicategory, not a cohomological closure phenomenon. We make the bicategorical
mismatch precise --- their distributive law is a \emph{mixed} comonad--monad law in
$[\Set,\Set]$ (an entwining/lifting), whereas $H^2$ classifies distributive laws of
two \emph{categories} in $\Span(\Set)$ --- and we delineate exactly the salvageable
positive content: among genuine \emph{directed-container-over-directed-container}
no-go theorems, the obstruction \emph{is} a class in $H^2(\Sk_\C;\D)$, realized by
the rigid twist as the generator of $\mathbb Z/2$. The correct grant statement is
the scoped one. All numerical claims are machine-checked
(\texttt{ks\_level\_check.py}, \texttt{cohomology\_holonomy.py}).
\end{abstract}

\section{The conjecture, and what would make it true}

The prove target read, verbatim in its success criteria:
\begin{enumerate}
\item Express list/stream/tree directed containers as small categories $\C$.
\item \emph{Express probability/nondeterminism monad as a wide subcategory $\D$.}
\item Verify or refute condition \textnormal{(L)}.
\item Compute the cocycle $\omega$ and show $[\omega]\neq0$ in $H^2(\Sk_\C;\D)$.
\end{enumerate}
Steps 3 and 4 presuppose step 2. The whole machinery of \cite{Gcohom,pairwise} ---
hom-presheaves $\Hom_\K(-,b)$ as right $\D$-sets, free bases, transversals, the
defect $2$-cochain, the complex $C^\bullet(\Sk_\C;\D)$ --- is \emph{defined relative
to a wide subcategory $\D\subseteq\K$}. If the monad cannot be a wide subcategory,
the conjecture does not typecheck, let alone hold. So the entire weight of the
conjecture rests on step 2. We show step 2 is impossible.

\paragraph{The only route to step 2.} A monad $T\colon\Set\to\Set$ ``is'' a category
in exactly one way relevant here: through the equivalence
\[
   \text{directed container}\ \simeq\ \text{small category}\ \simeq\
   \text{polynomial comonad}
\]
of Ahman--Chapman--Uustalu and Ahman--Uustalu \cite{acu,auCats}. A small category
$\D$ presents the polynomial endofunctor $X\mapsto\coprod_{a\in\ob\D}X^{\,\D(a,-)}$
(``shapes $=$ objects, positions $=$ outgoing arrows''), with its canonical comonad
structure. \emph{This is the sense in which the right factor $\D$ of a Zappa--Sz\'ep
product is an endofunctor of $\Set$.} For the Karamlou--Shah question ``does the
directed container distribute over $T$'' to become an instance of a strict
factorization, the monad $T$ ($=\Dist$ or $\Pow$) must be realized as such a
$\D$ --- i.e.\ $T$ must be a polynomial functor. (The orientation is irrelevant:
whichever of $\C,\D$ one assigns to $T$, that factor must be a small category, hence
$T$ must be polynomial. See Remark~\ref{rem:orient}.)

We now prove the route is blocked.

\section{$\Pow$ and $\Dist$ are not polynomial functors}

Recall a functor $F\colon\Set\to\Set$ is \emph{polynomial} (a \emph{container}) iff
it is a coproduct of representables, $F(X)\cong\coprod_{s\in S}X^{P_s}$ for a set $S$
of shapes and position-sets $P_s$; equivalently iff $F$ preserves connected limits
(wide pullbacks) \cite{aag,gambinokock}. Two elementary invariants suffice here:
$F(1)\cong S$ and $F(0)\cong\{s:P_s=\varnothing\}$.

\begin{proposition}[Powerset is not polynomial]\label{prop:Pow}
The finite (or full) covariant powerset functor $\Pow$ is not polynomial.
\end{proposition}

\begin{proof}
Suppose $\Pow\cong\coprod_{s\in S}(-)^{P_s}$. Evaluating: $\Pow(1)=2$ forces
$|S|=2$. $\Pow(0)=1$ (the only subset of $\varnothing$ is $\varnothing$) forces
\emph{exactly one} shape to have an empty position-set; call the other shape's
position-set size $k\ge1$. Then
\[
   \Pow(2)\ =\ 2^{0}+2^{k}\ =\ 1+2^{k}.
\]
But $\Pow(2)=4$ (the four subsets of a two-element set), so $2^{k}=3$, which has no
solution in integers $k$. Contradiction.
\end{proof}

\begin{proposition}[Distribution is not polynomial]\label{prop:Dist}
The finitely-supported distribution monad $\Dist$ is not polynomial.
\end{proposition}

\begin{proof}
$\Dist(1)=1$ (the unique distribution on a point) forces $|S|=1$ if $\Dist$ were
polynomial, i.e.\ $\Dist\cong(-)^{P}$ would be \emph{representable}. A representable
functor preserves all limits, in particular finite products:
$(-)^P$ sends $A\times B$ to $A^P\times B^P=(A^P)\times(B^P)$, and the comparison
$(A\times B)^P\to A^P\times B^P$ is the identity (an iso). We exhibit that the
corresponding comparison for $\Dist$ --- the marginal map
$\langle\Dist\pi_A,\Dist\pi_B\rangle\colon\Dist(A\times B)\to\Dist(A)\times\Dist(B)$
--- is not injective for $A=B=2$. Take the two joint distributions on
$\{0,1\}\times\{0,1\}$
\[
   \mu_{\mathrm{ind}}=\tfrac14\delta_{(0,0)}+\tfrac14\delta_{(0,1)}
     +\tfrac14\delta_{(1,0)}+\tfrac14\delta_{(1,1)},\qquad
   \mu_{\mathrm{cor}}=\tfrac12\delta_{(0,0)}+\tfrac12\delta_{(1,1)} .
\]
Both have uniform marginals $(\tfrac12,\tfrac12)$ on each factor, yet
$\mu_{\mathrm{ind}}\neq\mu_{\mathrm{cor}}$. So the marginal map is non-injective,
$\Dist$ does not preserve the product $2\times2$, hence is not representable, hence
(being $1$ on a singleton) not polynomial.
\end{proof}

Both computations are confirmed in \texttt{ks\_level\_check.py}: for $\Pow$,
$1+2^k=4\Rightarrow 2^k=3$ has no integer solution; for $\Dist$, the two displayed
joints have equal marginals and are distinct.

\begin{remark}[Why this is the heart of the matter]
$\Pow$ and $\Dist$ are the two canonical \emph{non-polynomial} $\Set$-monads ---
$\Pow$ fails to preserve wide pullbacks, $\Dist$ fails to preserve products. The
entire directed-container / $\Poly$ / $\mathbf{Cat}^{\#}$ world in which our
$H^2$ obstruction is defined is, by construction, the world of \emph{familially
representable} structure \cite{shapirospivak}. The Karamlou--Shah monads are
precisely the standard inhabitants of its complement. The conjecture asked for the
obstruction in a world the obstruction cannot see.
\end{remark}

\section{The level theorem}

\begin{theorem}[Level theorem]\label{thm:level}
Let the right factor of a sought strict factorization be a small category $\D$, as
required by the criterion of \cite{pairwise,Gcohom}. Then:
\begin{enumerate}
\item \textnormal{(Negative.)} The Karamlou--Shah right factors $\Dist$ and $\Pow$
cannot be presented as such a $\D$: by Propositions~\ref{prop:Pow}--\ref{prop:Dist}
they are not polynomial functors, hence not directed containers, hence not small
categories acting by precomposition on $\Set$. Consequently success-criterion~$2$
is unsatisfiable; conditions \textnormal{(L)} and \textnormal{(G)} are not even
\emph{defined} for these monads, and the proposed identification
``\,$\text{K--S no-go}=\text{(G) failure}=[\omega]\neq0$\,'' does not typecheck. The
Karamlou--Shah obstruction lies strictly \emph{below} \textnormal{(L)}: it is a
\emph{representability} obstruction in the bicategory $[\Set,\Set]$, prior to any
cohomological closure datum.
\item \textnormal{(Positive.)} When both factors \emph{are} directed containers
(small categories), the existence question for a complementary $\C$ with
$\K=\C\bowtie\D$ is governed, under hypothesis \textnormal{(H)} of \cite{Gcohom}, by
$[\omega]\in H^2(\Sk_\C;\D)$, and there are genuine \emph{no-go} instances with
$[\omega]\neq0$. The rigid twist of \cite[Ex.~5.1]{Gcohom} realizes the generator of
$H^2(\Sk_\C;\D)\cong\mathbb Z/2$.
\end{enumerate}
\end{theorem}

\begin{proof}
(1) Immediate from Propositions~\ref{prop:Pow}--\ref{prop:Dist} together with the
equivalence directed container $\simeq$ small category $\simeq$ polynomial comonad
\cite{acu,auCats}: a wide subcategory $\D\subseteq\K$ is in particular a small
category, whose induced $\Set$-endofunctor is polynomial; $\Dist,\Pow$ are not
polynomial, so no wide subcategory induces them. The conditions \textnormal{(L)},
\textnormal{(G)} are predicates on the pair $(\K,\D)$ with $\D$ a wide subcategory,
so they have no value when $\D$ does not exist.

(2) Is Theorem~5.2 (the rigid twist) of \cite{Gcohom}, recalled in
\S\ref{sec:positive}: there $\D=\E$ is the endomorphism subcategory, a directed
container ($\mathbb Z/2$ at $a$, trivial elsewhere); \textnormal{(L)} holds;
\textnormal{(G)} fails; $H^2(\Sk_\C;\D)\cong\mathbb Z/2$ and the defect represents
its nonzero element. Both factors $\C,\D$ are small categories, hence directed
containers.
\end{proof}

\begin{remark}[Orientation independence]\label{rem:orient}
One might hope to dodge Theorem~\ref{thm:level}(1) by assigning the monad $T$ to the
\emph{left} factor $\C$ instead of the right factor $\D$. This does not help: in a
strict factorization $\K=\C\bowtie\D$ \emph{both} $\C$ and $\D$ are wide
subcategories of $\K$, hence both are small categories, hence both induce polynomial
endofunctors. Realizing $T=\Dist$ or $\Pow$ as either factor requires $T$ polynomial.
The obstruction of Propositions~\ref{prop:Pow}--\ref{prop:Dist} is symmetric in the
two factors.
\end{remark}

\section{The bicategorical mismatch}

The deepest way to see why the conjecture mis-fires is that it silently translates
between two different notions of distributive law living in two different
bicategories.

\paragraph{(A) The Zappa--Sz\'ep law (where $H^2$ lives).} A strict factorization
$\K=\C\bowtie\D$ is the same datum as a distributive law of \emph{categories}
$\lambda\colon\D\C\Rightarrow\C\D$ --- a matched pair \cite[\S4]{pairwise}. Regarding
a small category as a monad in the bicategory $\Span(\Set)$ on its object set, this
is a distributive law of \emph{two monads in $\Span(\Set)$}, and the classical
correspondence ``distributive laws $\leftrightarrow$ strict factorizations''
(Rosebrugh--Wood \cite{rw}, Street \cite{street}) applies. Both participants are
monads (categories). The obstruction $[\omega]\in H^2(\Sk_\C;\D)$ measures the
failure of a transversal of $\C$-generators to close.

\paragraph{(B) The Karamlou--Shah law (what they actually study).} A directed
container is a \emph{comonad} $G$ on $\Set$. ``$G$ distributes over the monad $T$''
is a \emph{mixed} distributive law $\lambda\colon GT\Rightarrow TG$ (or its
companion $TG\Rightarrow GT$) in the endofunctor bicategory $[\Set,\Set]$ --- a
comonad and a monad, not two monads. Such a law is exactly an
\emph{entwining}/\emph{lifting} datum (Beck \cite{beck}; Brze\-zi\'nski--Majid
\cite{bm}; Power--Watanabe \cite{pw}): it lifts $T$ to the category of
$G$-coalgebras and lifts $G$ to the category of $T$-algebras, and its (co)algebras
are \emph{bialgebras}. It does \emph{not} present any category $\K$ as a factorized
$\C\bowtie\D$; the composite $GT$ is in general neither a monad nor a comonad. So
(B) is not an object that the $H^2$ theorem of (A) classifies, \emph{independently}
of whether $T$ is polynomial.
\smallskip

Thus there are two independent reasons the conjecture fails, and they reinforce each
other: \emph{(i)} $T=\Dist,\Pow$ is not even a category, so cannot be a factor
(Theorem~\ref{thm:level}(1)); and \emph{(ii)} even a perfectly good comonad--monad
law is the wrong \emph{kind} of distributive law --- a lifting, not a factorization.
The cohomological obstruction answers a question (existence of a complementary wide
subcategory) that the Karamlou--Shah setup never poses.

\begin{remark}[Two near-misses, ruled out]
\emph{(Kleisli/Lawvere presentation.)} The one way to make $T$ ``a category'' other
than as a directed container is to take its Kleisli category $\mathrm{Kl}(T)$ (so
$\mathrm{Kl}(\Pow)=\mathbf{Rel}$, $\mathrm{Kl}(\Dist)=$ stochastic maps) or its
Lawvere theory. Neither yields a Zappa--Sz\'ep instance. A wide subcategory must
share the object set of $\K$; $\mathrm{Kl}(T)$ has \emph{all sets} as objects and is
not small, while a directed container $\C$ has the shape set (e.g.\ $\mathbb N$ for
stream/list) as objects --- there is no common $\K$ in which both are wide.
Moreover the question ``does $\C$ admit a law over $\mathrm{Kl}(T)$'' is again a
lifting question (case (B)), not a factorization. So this route reproduces both
obstructions rather than escaping them.
\end{remark}

\section{What survives: directed-container no-go theorems \emph{are} $H^2$}
\label{sec:positive}

The positive half of Theorem~\ref{thm:level} is the grant-relevant residue, and it
is genuinely a ``no-go theorem with a computable cohomological obstruction'' --- only
for the right class of monads.

\begin{example}[A directed container that does not distribute, with obstruction
$\mathbb Z/2$]\label{ex:rt}
Objects $a,x,y$; $\End(a)=\{\id_a,g\}$ with $g^2=\id_a$, other endomorphism monoids
trivial; $\Hom(a,x)=\{p,pg\}$, $\Hom(a,y)=\{q,qg\}$, $\Hom(x,y)=\{s,s_2\}$ with the
\emph{rigid twist} $s\circ p=q,\ s_2\circ p=qg$. Take $\D=\E$, the endomorphism
subcategory: a directed container, equal to $\mathbb Z/2$ concentrated at $a$. Then
\textnormal{(L)} holds (every hom-set is free over its source vertex group),
hypothesis \textnormal{(H)} holds (abelian vertex groups, trivial left action), and
\[
   H^2(\Sk_\C;\D)\cong\mathbb Z/2,\qquad [\omega]\neq0,
\]
so no complementary directed container $\C$ makes $\K=\C\bowtie\D$: the directed
container $\E$ admits no Zappa--Sz\'ep complement inside $\K$. The four valid
transversals all yield a defect outside $B^2$ representing the same nonzero class
(machine-checked, \texttt{cohomology\_holonomy.py}: $\dim_{\mathbb F_2}H^2=1$, no
closing transversal).
\end{example}

This is the honest analogue, \emph{inside the polynomial world}, of a
Karamlou--Shah statement: a directed container (here $\E$) for which a desired
distributive law (here the matched pair assembling $\K$) provably does not exist,
\emph{and} whose obstruction is an explicit, computable cohomology class. The
contrast with Karamlou--Shah is exactly the content of Theorem~\ref{thm:level}: when
the second factor is itself a directed container, the obstruction is the $H^2$ class;
when it is $\Dist$ or $\Pow$, there is no second factor at all and the failure is one
level deeper.

\section{Honest scope}

We have \emph{not} reproved the Karamlou--Shah no-go theorems (that no mixed
distributive law $GT\Rightarrow TG$ exists for $G$ the list/stream/tree container
and $T=\Dist,\Pow$); that is a theorem about liftings in $[\Set,\Set]$, established
by them and not reconsidered here (we did not have access to their proof, and worked
only from the published statement: the monads are $\Dist$ and $\Pow$). What we have
established is sharper than a reproof: the \emph{type} of their obstruction. It is
\emph{not} the cohomological closure class $[\omega]\in H^2(\Sk_\C;\D)$; it cannot
be, because their second factor is not a directed container and their law is not a
factorization. The two facts we contribute --- the elementary non-polynomiality of
$\Pow,\Dist$ (Props.~\ref{prop:Pow}--\ref{prop:Dist}) and the bicategorical
distinction between mixed laws and factorizations (\S4) --- are each independently
sufficient to refute the conjectured identification, and both are machine- or
hand-checkable.

\paragraph{Consequence for the grant narrative.} The slogan ``compositional
correctness has a computable obstruction'' stands, correctly scoped: \emph{for
directed-container distributive laws --- the polynomial / $\mathbf{Cat}^{\#}$
world --- the obstruction to a strict factorization is the computable class
$[\omega]\in H^2(\Sk_\C;\D)$, vanishing iff a complement exists.} The probability
and nondeterminism no-go theorems are a \emph{different} phenomenon, a
non-representability obstruction outside this world, and conflating the two would
have been an overclaim. Locating the boundary precisely is the result.

\begin{thebibliography}{99}
\bibitem{Gcohom} MacBeth, \emph{The global-closure obstruction (G) is a cohomology
class: the Zappa--Sz\'ep holonomy as $H^2(\Sk_\C;\D)$}, Kodamai note, 11 June 2026.
\bibitem{pairwise} MacBeth, \emph{The pairwise Zappa--Sz\'ep criterion}, Kodamai
note, 10 June 2026.
\bibitem{acu} D.~Ahman, J.~Chapman, T.~Uustalu, \emph{When is a container a
comonad?}, Logical Methods in Computer Science 10(3) (2014).
\bibitem{auCats} D.~Ahman, T.~Uustalu, \emph{Directed containers as categories},
MSFP 2016, EPTCS 207, 89--98.
\bibitem{aag} M.~Abbott, T.~Altenkirch, N.~Ghani, \emph{Containers: constructing
strictly positive types}, Theoret.\ Comput.\ Sci.\ 342 (2005), 3--27.
\bibitem{gambinokock} N.~Gambino, J.~Kock, \emph{Polynomial functors and polynomial
monads}, Math.\ Proc.\ Cambridge Philos.\ Soc.\ 154 (2013), 153--192.
\bibitem{shapirospivak} D.~Spivak, \emph{Polynomial functors: a mathematical theory
of interaction} (and B.~Shapiro--D.~Spivak, \emph{Duoidal structures \dots/
$\mathbf{Cat}^{\#}$}), 2023--2024.
\bibitem{rw} R.~Rosebrugh, R.~J.~Wood, \emph{Distributive laws and factorization},
J.\ Pure Appl.\ Algebra 175 (2002), 327--353.
\bibitem{street} R.~Street, \emph{The formal theory of monads}, J.\ Pure Appl.\
Algebra 2 (1972), 149--168.
\bibitem{beck} J.~Beck, \emph{Distributive laws}, Lecture Notes in Math.\ 80 (1969),
119--140.
\bibitem{bm} T.~Brze\-zi\'nski, S.~Majid, \emph{Coalgebra bundles}, Comm.\ Math.\
Phys.\ 191 (1998), 467--492.
\bibitem{pw} J.~Power, H.~Watanabe, \emph{Combining a monad and a comonad}, Theoret.\
Comput.\ Sci.\ 280 (2002), 137--162.
\bibitem{ks} P.~Karamlou, A.~Shah, \emph{No-go theorems: directed containers that
do not distribute over distribution monads}, LICS 2024.
\end{thebibliography}

\end{document}
